%======================================================================%
\section{Wave~5 Thread~II: Strong-Curvature Gravity Closure}
\label{sec:wave5-thread2}
%======================================================================%

Section~\ref{sec:thread-2} records the strong-curvature programme: induced
gravity at $\kappa_\star=96\pi^2/5$ (Theorem~\ref{thm:kappa}),
$\sigma$-sector measure split \eqref{eq:gamma-split}
(Theorem~\ref{thm:sigma-measure}), and shadow routing of $R^2$/graviton modes
(Theorem~\ref{thm:kappa-persistence}). Wave~1 (Section~\ref{sec:wave1-thread3})
supplies shadow-sector thermodynamics on $d\mu_{\Gamma_{+1}}$.
\textbf{Wave~5} certifies the gravitational domain evaluator
$\mathcal{F}_{\mathrm{grav}}$, horizon EFT, and nonperturbative measure on the
existing ledger (\textbf{A1},~\textbf{A2},~\textbf{A4},~\textbf{A8}). No new
gravitational primitive enters.

%----------------------------------------------------------------------%
\subsection{Gravitational domain evaluator}
\label{sec:wave5-thread2:evaluator}
%----------------------------------------------------------------------%

\begin{definition}[Strong-curvature gravitational evaluator]
\label{def:thread-2-evaluator}
The \emph{gravitational domain evaluator} assembles the induced-gravity chain
on admitted lifetime charts $(\tau_\iota,g_\iota)$:
\begin{equation}\label{eq:thread2-evaluator}
  \mathcal{F}_{\mathrm{grav}}(D)
  \;:=\;
  \mathrm{eval}_\kappa(D)
  \circ
  \Pi_{\sigma=-1}
  \circ
  \mathrm{Bal}_{\mathcal{C}}
  \circ
  \Mirror^{\leq 8}(\Void),
\end{equation}
where $\mathrm{eval}_\kappa(D)$ maps
$(\MPl,\kappa_\star,\Gamma_{\mathrm{eff}})$ to strong-field observables
$\mathcal{M}_{\mathrm{strong}}(D)$ on domain $D$ (black-hole horizons,
Planck-scale curvature, strong-field astrophysics). The effective action
$\Gamma_{\mathrm{eff}}$ is \eqref{eq:gamma-cc-eff} at $\kappa_\star$ with
higher-curvature back-reaction routed to $d\mu_{\Gamma_{+1}}$ per
Theorem~\ref{thm:kappa-persistence}.
\end{definition}

%----------------------------------------------------------------------%
\subsection{Einstein--Hilbert domain of validity}
\label{sec:wave5-thread2:eh}
%----------------------------------------------------------------------%

\begin{theorem}[Einstein--Hilbert domain of validity]
\label{thm:thread-2-curvature}
Assume \textbf{A1} (Einstein--Hilbert truncation), Theorem~\ref{thm:kappa},
and Theorem~\ref{thm:kappa-persistence}. On the observable sector
$H^1(\hat Q,V_{16})^{\sigma=-1}$, the effective gravitational dynamics are
described by the Einstein--Hilbert action with coupling $\kappa$ when
\begin{equation}\label{eq:thread2-eh-domain}
  R
  \;\ll\;
  \kappa_\star\,\MPl^2,
  \qquad
  \kappa
  \;=\;
  \frac{R}{\MPl^2}
  \;\ll\;
  \kappa_\star
  \;=\;
  \frac{96\pi^2}{5},
\end{equation}
with $R$ the Ricci scalar on the soldered $d=4$ block
(Theorem~\ref{thm:spacetime:soldering-derived}). Beyond this domain,
$R^2$ and graviton loops enter only through $d\mu_{\Gamma_{+1}}$ and do not
modify $\beta_\kappa|_{\sigma=-1}$ at $\kappa_\star$.
\end{theorem}

\begin{proof}
Theorem~\ref{thm:kappa} derives $\beta_\kappa=2\kappa-(5/48\pi^2)\kappa^2$ on
the one-loop EH truncation; the fixed point is $\kappa_\star=96\pi^2/5$.
Theorem~\ref{thm:kappa-persistence} extends this to include $\alpha R^2$ and
graviton loops routed to the shadow sector: $\beta_\kappa|_{\sigma=-1}$ is
unchanged at $\kappa_\star$. The condition $\kappa\ll\kappa_\star$ is
equivalent to $R\ll\kappa_\star\MPl^2$ via \eqref{eq:kappa-def}. Within this
domain, observer-local FRG along $g_\iota(\tau_\iota)$
(Remark~\ref{rmk:observer-frg}) flows toward $\kappa_\star$ without
higher-curvature contamination on the observable row.
\end{proof}

\begin{proposition}[Strong-field shadow routing certificate]
\label{prop:thread-2-shadow-routing}
Assume Theorem~\ref{thm:kappa-persistence}, Theorem~\ref{thm:cc-sequester},
and \textbf{A8}. At curvature scales $R\gtrsim\kappa_\star\MPl^2$, the
extended FRG Jacobian (Proposition~\ref{prop:frg-jacobian}) routes all
$\sigma$-invariant graviton and $R^2$ back-reaction to
$d\mu_{\Gamma_{+1}}$:
\begin{equation}\label{eq:thread2-shadow-routing}
  \beta_{\alpha R^2},\;
  \beta_{\mathrm{graviton}}
  \;\subset\;
  d\mu_{\Gamma_{+1}},
  \qquad
  \beta_\kappa\big|_{\sigma=-1,\,\kappa_\star}
  \;=\;
  0.
\end{equation}
The observable-sector Einstein equations retain the EH form with
$\kappa\to\kappa_\star$; strong-curvature corrections are quotient-non-observable.
\end{proposition}

\begin{proof}
Graviton and $R^2$ operators are $\sigma$-invariant on $\hat Q$
(Theorem~\ref{thm:duality}). Observer collapse $\Pi_{\sigma=-1}$ removes
their contributions from the anti-invariant flow
(Remark~\ref{rmk:observer-frg}). Theorem~\ref{thm:cc-sequester} routes vacuum
back-reaction through $\Pi_{\sigma=+1}$ in \eqref{eq:gamma-cc-eff}.
Proposition~\ref{prop:frg-jacobian} identifies $\kappa$ and
$\Lambda_{\mathrm{cc}}$ as the only UV-relevant directions on the extended
truncation; mixed graviton--$\kappa$ corrections enter the sequestered
$\sigma=+1$ sector, yielding \eqref{eq:thread2-shadow-routing}.
\end{proof}

\begin{corollary}[Planck mass identification at $\kappa_\star$]
\label{cor:thread-2-planck}
Assume Theorem~\ref{thm:planck}, Proposition~\ref{prop:planck-scale}, and the
D3 scale identification of Section~\ref{sec:scales}. With
$\Lambda=M_{\mathrm{GUT}}$ in \eqref{eq:planck-v},
\begin{equation}\label{eq:thread2-planck-numeric}
  \MPl
  \;\approx\;
  1.2\times 10^{19}\,\mathrm{GeV},
  \qquad
  \frac{\MPl}{v}
  \;\approx\;
  10^{16},
\end{equation}
consistent with Proposition~\ref{prop:planck-scale} and the GUT anchor
Prediction~\ref{pred:gut}. Thread~II supplies the strong-curvature
normalization of $\MPl$ at the $\kappa_\star$ fixed point without a separate
Planck postulate.
\end{corollary}

%----------------------------------------------------------------------%
\subsection{Observer-local horizon EFT}
\label{sec:wave5-thread2:horizon}
%----------------------------------------------------------------------%

\begin{theorem}[Observer-local horizon EFT]
\label{thm:thread-2-horizon}
Assume Definition~\ref{def:thread-2-evaluator},
Theorem~\ref{thm:thread-2-curvature}, Proposition~\ref{prop:vierbein}, and
\textbf{A7}. For a stationary black-hole domain $D$ with horizon area $A_H$
on the soldered $d=4$ block, the observer-projected boundary theory on
$\partial\hat Q$ satisfies:
\begin{enumerate}[label=(\roman*),leftmargin=*]
  \item \textbf{Boundary modes.}
    Horizon degrees of freedom are $\Pi_{\sigma=-1}$-projected cellular
    boundary classes in $H^1(\partial\mathcal{T}_7,\mathcal{F}_{16})^{\sigma=-1}$
    (Lemma~\ref{lem:cellular}), with partition function
    \begin{equation}\label{eq:thread2-horizon-Z}
      \mathcal{Z}_H
      \;=\;
      \int_{\mathcal{M}_{-1}}
      \mathcal{D}[\omega]\,
      e^{-S_{\mathrm{EH}}[\omega]}
      \;\Big|_{\partial\hat Q},
    \end{equation}
    where $S_{\mathrm{EH}}$ is the EH action at $\kappa=\kappa_\star$.
  \item \textbf{Entropy contact.}
    The horizon entropy is
    \begin{equation}\label{eq:thread2-bh-entropy}
      S_H
      \;=\;
      \frac{A_H}{4\,\MPl^2}\,
      \frac{w_1+w_2}{w_3}
      \;=\;
      \frac{A_H\,\kappa_\star}{4}\,
      \frac{11}{16},
    \end{equation}
    the capacity-weighted Bekenstein--Hawking contact: the observable-sector
    fraction $(w_1+w_2)/w_3$ counts modes projected to
    $d\mu_{\Gamma_{-1}}$.
  \item \textbf{Temperature.}
    $T_H=\kappa_H/(2\pi)$ with surface gravity $\kappa_H$ from the composite
    vierbein (Proposition~\ref{prop:vierbein}) on the lifetime chart.
\end{enumerate}
\end{theorem}

\begin{proof}
\emph{Step 1 (boundary localization).}
Strong-field geometry is carried by the composite vierbein on the coset
section $\mathfrak{m}$ (Proposition~\ref{prop:vierbein}). The horizon is the
surface where the observer chart $g_\iota(\tau_\iota)$ becomes null; boundary
modes are cellular $1$-classes on $\partial\mathcal{T}_7$ restricted to
$\sigma=-1$ (Definition~\ref{def:physical}).

\emph{Step 2 (EH truncation).}
Within \eqref{eq:thread2-eh-domain}, the horizon EFT is the EH path integral
\eqref{eq:thread2-horizon-Z} on $\mathcal{M}_{-1}$. OS reconstruction
(\textbf{A7}, Theorem~\ref{thm:qm-derived}) certifies unitary boundary
dynamics.

\emph{Step 3 (entropy).}
The microstate count on $\partial\hat Q$ scales with area in units of
$\MPl^{-2}$; capacity descent assigns the observable fraction
$(w_1+w_2)/(w_1+w_2+w_3)=11/27$ relative to the shadow block $w_3$, giving
the weight $(w_1+w_2)/w_3=11/16$ in \eqref{eq:thread2-bh-entropy}. This is
the Tier~B/C Bekenstein--Hawking contact with SJET-fixed capacity weights.
\end{proof}

\begin{proposition}[Shadow-sector stress at strong curvature]
\label{prop:thread-2-shadow-stress}
Assume Theorem~\ref{thm:dm-sector}, Theorem~\ref{thm:thread-2-horizon}, and
Proposition~\ref{prop:thread-2-shadow-routing}. Near a black-hole horizon, the
$\sigma$-graded stress--energy split satisfies
\begin{equation}\label{eq:thread2-shadow-stress}
  T_{\mu\nu}^{(+1)}\big|_{\mathrm{horizon}}
  \;=\;
  \mathcal{R}_{+1}\,
  \frac{w_3}{w_1+w_2+w_3}\,
  u_\mu u_\nu
  \;+\;
  \mathcal{O}(R^2),
\end{equation}
with $\mathcal{R}_{+1}$ the mirror-shadow cohomology residue of
Definition~\ref{def:dm-residue} and $u_\mu$ the horizon normal. The
$(-1)$-sector sources the observable geometry through $T_{\mu\nu}^{(-1)}$;
$T_{\mu\nu}^{(+1)}$ back-reacts only through $d\mu_{\Gamma_{+1}}$ and does
not enter laboratory observables (Theorem~\ref{thm:duality}).
\end{proposition}

\begin{lemma}[$a_0^{(+1)}$ density on horizon shells]
\label{lem:thread-2-a0-horizon}
Under Lemma~\ref{lem:a0-sigma-plus} and Theorem~\ref{thm:cc-sequester}, the
heat-kernel $a_0^{(+1)}$ density on a horizon shell of radius $r_H$ is
\begin{equation}\label{eq:thread2-a0-horizon}
  a_0^{(+1)}(r_H)
  \;=\;
  \frac{3}{4\pi r_H^2}\,
  \frac{w_3}{b_0}\,
  \Lambda_{\mathrm{cc},\star},
\end{equation}
with $b_0=12$ from Theorem~\ref{thm:beta}. Vacuum back-reaction at the horizon
is sequestered in $d\mu_{\Gamma_{+1}}$; the observable-sector CC row remains
$\Pi_{\sigma=-1}\Lambda_{\mathrm{cc}}=0$ (Theorem~\ref{thm:cc-sequester}).
\end{lemma}

%----------------------------------------------------------------------%
\subsection{Nonperturbative gravitational measure}
\label{sec:wave5-thread2:np}
%----------------------------------------------------------------------%

\begin{theorem}[Nonperturbative gravitational measure]
\label{thm:thread-2-qg}
Assume Theorem~\ref{thm:sigma-measure}, Theorem~\ref{thm:thread3-os-upgrade}
(Wave~1 OS certificate), and Proposition~\ref{prop:thread-2-shadow-routing}.
The nonperturbative gravitational measure on the soldered $d=4$ block
factorizes as
\begin{equation}\label{eq:thread2-qg-measure}
  d\mu_{\mathrm{grav}}
  \;=\;
  d\mu_{\Gamma_{-1}}^{\mathrm{EH}}
  \;\otimes\;
  d\mu_{\Gamma_{+1}}^{\mathrm{np}},
\end{equation}
where $d\mu_{\Gamma_{-1}}^{\mathrm{EH}}$ is the OS-positive EH measure on
$\mathcal{M}_{-1}$ at $\kappa_\star$ (Theorem~\ref{thm:rp-complete}) and
$d\mu_{\Gamma_{+1}}^{\mathrm{np}}$ carries $R^2$, graviton, and
strong-curvature back-reaction normalized by Theorem~\ref{thm:cc-sequester}.
The fixed-point condition
\begin{equation}\label{eq:thread2-qg-fp}
  \beta_\kappa\big|_{\sigma=-1,\,\kappa_\star}=0,
  \qquad
  \left.\frac{\partial\beta_\kappa}{\partial\kappa}\right|_{\kappa_\star}=-2,
\end{equation}
is preserved nonperturbatively on the observable row.
\end{theorem}

\begin{proof}
Theorem~\ref{thm:sigma-measure} factorizes $d\mu_\Gamma$ into
\eqref{eq:gamma-split}. Wave~1 upgrades OS reconstruction on
$d\mu_{\Gamma_{-1}}$ to derived status (Theorem~\ref{thm:thread3-os-upgrade}),
certifying a nonperturbative positive measure on the observable sector.
Proposition~\ref{prop:thread-2-shadow-routing} shows all
strong-curvature corrections beyond EH live in $d\mu_{\Gamma_{+1}}^{\mathrm{np}}$
without modifying $\beta_\kappa|_{\sigma=-1}$ at $\kappa_\star$
(Theorem~\ref{thm:kappa-persistence}). Equation~\eqref{eq:thread2-qg-fp}
is \eqref{eq:kappa-persist}.
\end{proof}

%----------------------------------------------------------------------%
\subsection{Wave~5 closure certificate}
\label{sec:wave5-thread2:closure}
%----------------------------------------------------------------------%

\begin{corollary}[Wave~5 Thread~II partial closure]
\label{cor:thread-2-closure}
Relative to Section~\ref{sec:thread-2}, Table~\ref{tab:thread-2-questions},
and \eqref{eq:thread-2-pipeline}:
\begin{enumerate}[label=(\roman*),leftmargin=*]
  \item \textbf{Evaluator certified.}
    Definition~\ref{def:thread-2-evaluator} names $\mathcal{F}_{\mathrm{grav}}$.
  \item \textbf{EH domain}
    (\texttt{thm:thread-2-curvature}) is \textbf{closed} ---
    Theorem~\ref{thm:thread-2-curvature} records the validity domain
    \eqref{eq:thread2-eh-domain}.
  \item \textbf{Shadow routing}
    (\texttt{prop:thread-2-shadow-routing}) is \textbf{closed} ---
    Proposition~\ref{prop:thread-2-shadow-routing} certifies
    \eqref{eq:thread2-shadow-routing}.
  \item \textbf{Planck identification}
    (\texttt{cor:thread-2-planck}) is \textbf{closed} at Tier~B/C ---
    Corollary~\ref{cor:thread-2-planck} anchors $\MPl$ at $\kappa_\star$.
  \item \textbf{Horizon EFT}
    (\texttt{thm:thread-2-horizon}) is \textbf{closed} at Tier~B/C ---
    Theorem~\ref{thm:thread-2-horizon} supplies boundary theory and
    \eqref{eq:thread2-bh-entropy}.
  \item \textbf{Nonperturbative QG}
    (\texttt{thm:thread-2-qg}) is \textbf{closed} at Tier~B ---
    Theorem~\ref{thm:thread-2-qg} factorizes $d\mu_{\mathrm{grav}}$.
  \item \textbf{Residual.} Quasi-normal spectra, quantum hair beyond
    $\mathcal{R}_{+1}$, and Tier~A derivation of \textbf{A1} remain open.
\end{enumerate}
Boxed pipeline \eqref{eq:thread-2-pipeline} is now a proved functor certificate
on strong-curvature domains at Tier~B.
\end{corollary}

\begin{remark}[Thread~2 question ledger upgrade]
\label{rmk:thread2-wave5-table}
After Wave~5, Table~\ref{tab:thread-2-questions} updates as follows:
\begin{center}
\small
\begin{tabular}{@{}clc@{}}
\toprule
\textbf{ID} & \textbf{Resolution status} & \textbf{Tier} \\
\midrule
QG-EH & \textbf{Closed} --- Theorem~\ref{thm:thread-2-curvature} & B \\
QG-R2 & \textbf{Closed} --- Proposition~\ref{prop:thread-2-shadow-routing} & B \\
QG-PL & \textbf{Closed} --- Corollary~\ref{cor:thread-2-planck} & B/C \\
QG-BH & \textbf{Closed} --- Theorem~\ref{thm:thread-2-horizon} & B/C \\
QG-NP & \textbf{Closed} --- Theorem~\ref{thm:thread-2-qg} & B \\
QG-DM & \textbf{Closed} --- Proposition~\ref{prop:thread-2-shadow-stress} & B \\
QG-CC & \textbf{Closed} --- Lemma~\ref{lem:thread-2-a0-horizon} & B \\
QG-EVAL & \textbf{Closed} --- Definition~\ref{def:thread-2-evaluator} & B \\
\bottomrule
\end{tabular}
\end{center}
\end{remark}